\section{Attaching the funded service to the selected history}
\label{sec:original-history-portfolio-attachment}

The finite-volume service is useful only if its data belong to the history
selected for the public theorem. We verify that binding here. The public
construction selects one run and applies the funded continuation to that
run's own history, schedule, transitions, and finite-volume data. No cutoff,
period, or extraction order is used to select a second history. The common
schedule and its numerical prices are organized in
Section~\ref{sec:common-funded-history}; here we check the remaining
same-witness interface.

Let $R$ denote the selected run. Write $Y_R$ for its stored finite
renormalization-group history and $\mathcal P_R$ for the separately stored
scale campaign.
The distinction matters: the history used below is
$\operatorname{trail}(Y_R)$, whereas the campaign supplies the scale data.
Let $t_R(d)$ be the reference level at grade $d$, let $S_R(d,N,p)$ be the
transfer datum, and let $D_R(N,p,m)$ be the selected transition record. We
write
\[
 C_R=\bigl(\operatorname{trail}(Y_R),\ t_R,\ S_R\bigr)
\]
for the family with its fixed holonomy, label frames, and local sections.
This notation names the family already stored in the run; it does not define
a new history.

\begin{proposition}[Same-history attachment]
\label{prop:ex-original-history-attachment}
Suppose the selected run $R$ carries its funded continuation certificate.
Applying that certificate to the data stored in $R$ gives a package on the
same family $C_R$ with the following properties:
\begin{enumerate}
\item the reference-level condition
\[
 d+3\leq J(t_R(d))\qquad(d\geq0);
\]
\item one nesting family, one scale schedule, one transition family, and
   one transfer family, such that the selected transition at $(N,p,m)$ is
   available for $m<N$, and, for every $d,N,p$ with $t_R(d)\leq N$, the
   finite-volume transfer identity holds using those same data;
\item the funded nonlinear-response and adjacent-scale estimates, together
   with the selected service on exactly $C_R$;
\item the centered service, including its height estimates and terminal
   forced-profile relation, as well as the finite-volume bounds and marked
   portfolio conclusions recorded for $C_R$.
\end{enumerate}
The first three items and the finite-volume identities hold simultaneously
for all cutoffs and periods in their stated domains. For each fixed grade
and section in the purchased portfolio, the terminal marked estimate holds
for all sufficiently large cutoffs, uniformly in the period; its threshold may depend on that grade
and section. It is a separate local estimate, not a field included
implicitly in the centered service.
\end{proposition}

\begin{proof}
We first identify the history-bearing data. The selected run stores $Y_R$,
the schedule $t_R$, a transition family, the gradewise regularity data, and
the finite-volume transfer witnesses. Its underlying family is
\[
 C_R=\bigl(\operatorname{trail}(Y_R),t_R,S_R\bigr),
\]
with the fixed holonomy, label frames, and local-section map. The funded
continuation is a function of precisely these stored objects: the
configuration and its nesting data, the schedule and its opening
inequalities, the transfer family and transition records, and the
regularity and finite-volume transfer witnesses. Its output therefore has
the family formed from those same arguments. That family is $C_R$ itself;
there is no intervening identification map and no reselected label system.

Here is why the quantifiers in that application are uniform. The public
selection fixes the history $Y_R$ and a nesting family for it.
Independently, it chooses the reference schedule on that selected run so
that $d+3\leq J(t_R(d))$ for each grade. On this fixed history it constructs
the transition branches and their records once, and then defines one
transfer family on all triples $(d,N,p)$. That family has the required
regularity properties; for every $d,N,p$ with $t_R(d)\leq N$, its
finite-volume witness is built from the already chosen nesting family and
transition records. The witness at $(d,N,p)$ may depend on those indices,
but the history, nesting family, schedule, transitions, and transfer family
do not. Thus the simultaneous quantifiers are part of the selected run and
are not obtained by diagonal extraction or by making new global choices
after fixing $(N,p)$.

For the transfer witnesses, if $t_R(d)>N$ the construction uses a default
map. If $t_R(d)\leq N$, it composes along the selected transitions and
obtains the finite-volume identity and accounting estimate. This explains
the cutoff in item 2 and shows that both cases use the original $Y_R$.

The funded continuation supplies two quantitative inputs to the centered
construction: bounds on the delayed nonlinear response and its tariff, and
uniform estimates for each allowed adjacent-scale step. The latter are the
scale inequalities proved in
Section~\ref{ex-centered-geyser-scale-estimates}. That section applies them
to the selected adjacent levels and block size, and supplies a common
monotone threshold. It does not change the selected history, nor does it
by itself establish the nonlinear-response bounds.

The centered service is assembled from these data as follows. The escape
estimate uses the same nesting family, nonlinear-response bounds, and all
selected transitions. Its room and bank estimates use the budget fields of
the selected factor. The adjacent-scale inequalities give the required
scale comparison at each selected step. For heights at grade $d$ and
reference level $k$, the local residence estimate at $k$ gives the later
scale bounds, the one-cover footprint estimate, and a uniform camera
bound. The five numerical implications used in this conversion are proved
in Proposition~\ref{prop:ex-geyser-service-fields}; the eventual estimates
and finite-prefix maximum are in
Theorem~\ref{thm:ex-selected-geyser-steam}. The maximum is over the finitely
many grades below each target stage and its two adjacent stages, so it gives
a common threshold for all later uses at that stage.

At the terminal endpoint, the finite-volume transfer witness gives the top
transfer identity and its forced-profile relation. The centered service's
terminal estimate is this same relation, using the domination and rise
identities attached to the selected run. At an intermediate scale, the
transfer witnesses and selected transitions feed the local transition
estimate: specializing the tariff inequality to the zero-grade section
gives the required relay bounds. The separate growth and jump estimates
follow from positivity and the selected scale payment. The three-card
coefficient, exponential-entry, and forced-profile identities are proved
in Section~\ref{sec:three-port-history}; we use those identities here
rather than constructing a second interpolation or repeating their price
proof.

The funded continuation returns the centered service together with its
finite-volume bounds and marked portfolio. Its local transition estimates
use the same nonlinear-response bounds, adjacent-scale estimates, nesting
family, transition records, finite-volume witnesses, and schedule just
described. The core centered-service assembly and its stronger funded
version are distinguished in the source, but both use the supplied
configuration, schedule, transfer data, and transition family. Neither
creates a second history.

Finally, the terminal marked-field estimate is attached only after $C_R$
has been fixed. Section~\ref{sec:terminal-marked-construction} supplies the
reserved addresses and terminal representation from the regularity and
stock data on this same family. The quantitative local bounds and their
eventual uniformity for each fixed grade and section are proved in
Section~\ref{sec:ex-marked-field-local}. The local estimate requires soundness,
reference-level eligibility, fit of the union of the stencils, nonnegative
credit, a keeper, the regional card-times-allowance bound, cover at most
$81$, the square-root allowance-times-storm bound, and the stated positive
terminal weight and nonnegative reserve. The purchased local-estimates
section supplies these conditions for all sufficiently large cutoffs for
each fixed grade and section in the purchased portfolio, uniformly in the
period. The marked density
and forced-profile identity then give the terminal estimate on $C_R$.
Throughout, its history, schedule, transfer family, and terminal
representation remain those of $C_R$. No uniqueness assertion for a
terminal field is needed.
\end{proof}

\begin{remark}[Scope of the attachment]
The proposition proves the same-history attachment once the selected run's
funded continuation certificate is available. The common funded-history
construction is responsible for producing that certificate and its
quantitative prices; this section does not replace those purchases by the
name of a record. It also does not use the separate grantor route as
evidence for the public caller. The public construction applies its funded
continuation to the data stored in the selected run.
\end{remark}
